\chapter*{Introduction générale}
\addcontentsline{toc}{chapter}{Introduction générale}
\markboth{Introduction générale}{}

Dans un contexte de transformation numérique accélérée, les entreprises cherchent à
moderniser leurs processus de formation professionnelle. Les méthodes traditionnelles
~--- formations présentielles coûteuses, modules e-learning peu engageants~--- peinent à
répondre aux exigences d'efficacité et de standardisation des organisations modernes.
La Réalité Virtuelle (VR) émerge comme une réponse concrète à ces limites~: des études
empiriques indiquent que la formation en VR améliore la rétention des informations de
70~\% et rend les apprenants 4~fois plus rapides qu'une formation classique~\cite{pwc2020},
ce qui justifie son intégration dans les parcours de formation professionnelle B2B.

Le potentiel pédagogique de la VR étant établi, son adoption effective par les entreprises
se heurte toutefois à des obstacles concrets. Les solutions immersives disponibles sont
soit trop génériques~--- conçues indépendamment des procédures réelles d'une organisation
donnée~---, soit développées sur mesure pour un client unique, sans plateforme de gestion
permettant à l'entreprise de déployer la formation, d'en contrôler les accès et d'en
mesurer objectivement les résultats. Dans les deux cas, l'entreprise cliente reste sans
visibilité sur la progression réelle de ses collaborateurs, ce qui prive la formation de
la standardisation et de la traçabilité qu'elle devrait apporter. Le défi ne consiste
donc pas seulement à produire une expérience VR convaincante, mais à l'intégrer dans une
chaîne complète, de la gestion administrative des entreprises clientes jusqu'à la
restitution mesurable des performances de chaque employé~: comment concevoir une plateforme
B2B qui propose des formations professionnelles immersives en Réalité Virtuelle tout en
offrant aux entreprises les moyens de gérer leurs employés et de suivre objectivement leurs
performances~?

C'est à cette question que répond \tynassit{}, développée par Tynass~IT dans le cadre de
ce projet de fin d'études~: une plateforme B2B permettant aux entreprises clientes d'accéder
à des formations professionnelles immersives sur le casque Meta~Quest~3. Elle repose sur
trois composantes complémentaires. Un backend REST (Node.js, Express, MongoDB~Atlas) assure
l'authentification sécurisée, la gestion des accès et le suivi des sessions de formation.
Deux tableaux de bord web (React, TypeScript) offrent, l'un aux équipes Tynass la gestion
des entreprises et des modules VR, l'autre aux entreprises clientes la gestion de leurs
employés et la consultation des résultats. Enfin, une application VR déployée sur
Meta~Quest~3 permet aux employés d'effectuer des simulations immersives avec suivi
automatique des performances. Le premier module développé concerne l'intégration des
nouveaux employés d'Indigo~Company, leader tunisien du prêt-à-porter.

Ce rapport est organisé en six chapitres. Le premier présente le contexte général du
projet, une étude des solutions existantes, le choix de la méthodologie de travail~---
une approche Agile, itérative et incrémentale~--- ainsi que l'environnement de
développement. Le deuxième chapitre est consacré à la spécification des besoins et au
pilotage du projet. Les quatre chapitres suivants détaillent les quatre itérations de
développement réalisées, chacune apportant de nouvelles fonctionnalités à la plateforme~:
le backend et son architecture de sécurité, les tableaux de bord web, l'application VR de
base, puis la simulation d'intégration d'Indigo~Company. Une conclusion générale présente
enfin les résultats obtenus et les perspectives d'évolution du projet.

\clearpage